\documentclass[11pt,aspectratio=43,t]{beamer}

% -------------------------------------------------
% Theme and appearance
% -------------------------------------------------

\usetheme{Madrid}
\usecolortheme{default}
\setbeamertemplate{blocks}[rounded]
% Keep slides clean
\setbeamertemplate{navigation symbols}{}
\setbeamertemplate{footline}{}

% Wider margins for readability
\setbeamersize{text margin left=1.4cm, text margin right=1.4cm}

% Title hierarchy
\setbeamerfont{title}{series=\bfseries}
\setbeamerfont{frametitle}{series=\bfseries,size=\large}

% Bullet style and spacing (discourages overfull slides)
\setbeamertemplate{itemize items}[circle]
\setlength{\itemsep}{0.7em}
\setlength{\parskip}{0.3em}

% Do nothing automatically at new sections (we use explicit divider frames)
\AtBeginSection{}

% -------------------------------------------------
% Maths and graphics
% -------------------------------------------------
\usepackage{amsmath,amssymb}
\usepackage{graphicx}
\graphicspath{
  {figures/}
%  {figures/chap1/}
}
\usepackage{makecell}
\usepackage{booktabs}
% Optional: set your default figure path(s) here if you like
% \graphicspath{{figures/}{figures/lectures/}}



% -------------------------------------------------
% Title information (override per lecture file)
% -------------------------------------------------
\title[Geometry Without Manifolds]{What Becomes of Geometry\\When the Manifold Is Gone?}
%\subtitle{Combinatorial Mesh Calculus}
\author{\href{https://scholar.google.co.uk/citations?user=3nWJe5wAAAAJ&hl=en}{\textbf{Andrey Jivkov}}}
\institute{Department of Mechanical and Aerospace Engineering\\
The University of Manchester}
\date{\today}


% -------------------------------------------------
% Custom slide types / helpers
% -------------------------------------------------

% Title slide helper (so lecture files stay short and consistent)
\newcommand{\maketitleslide}{
\begin{frame}
  \titlepage
\end{frame}
}

% Overview slide helper (accepts an itemize body)
\newcommand{\overviewframe}[1]{
\begin{frame}{What this lecture is about}
\begin{itemize}
  #1
\end{itemize}
\end{frame}
}

% Section divider slide
\newcommand{\sectionframe}[1]{
\begin{frame}[plain]
  \vfill
  \centering
  {\Large\bfseries #1}
  \par\vspace{0.5em}
  \rule{0.4\textwidth}{0.4pt}
  \vfill
\end{frame}
}

% Figure placeholder slide
\newcommand{\figureframe}[1]{
\begin{frame}{}
  \centering
  \vspace{1.0cm}
  \textit{[Figure placeholder: #1]}
\end{frame}
}

% Figure + text frame (left text, right figure placeholder)
% #1 = frame title
% #2 = left content (typically itemize, keyidea, short text)
% #3 = figure description
\newcommand{\figtextframe}[3]{
\begin{frame}{#1}
\begin{columns}[T]
  \begin{column}{0.52\textwidth}
    #2
  \end{column}
  \begin{column}{0.46\textwidth}
    \centering
    \vspace{0.2em}
    \textit{[Figure: #3]}
  \end{column}
\end{columns}
\end{frame}
}

% Key idea emphasis (keep it short; do not paste full definitions)
\newenvironment{keyidea}{
\begin{itemize}
  \item[\textbf{Key idea}]
}{
\end{itemize}
}

% Summary slide helper
\newcommand{\takeawaysframe}[1]{
\begin{frame}{Key takeaways}
\begin{itemize}
  #1
\end{itemize}
\end{frame}
}

% -------------------------------------------------
% Document begins in the lecture file that inputs this template
% -------------------------------------------------